%! TEX root = ../edtc-18-19.tex

\chapter{Curbe în plan și în spațiu}


\section{Ecuații pentru curbe plane}

\index{curbe plane!ecuații parametrice}
\index{curbe plane!ecuații carteziene implicite}
\index{curbe plane!ecuații carteziene explicite}
\textbf{Curbele plane} pot fi reprezentate prin:
\begin{enumerate}[(a)]
  \item Ecuații parametrice, de forma:
    \[
      \alpha : \begin{cases}
        x &= x(t) \\
        y &= y(t)
      \end{cases}, t \in (a, b)
    \]
  \item Ecuații carteziene implicite, de forma $ y = f(x) $;
  \item Ecuație carteziană explicită, de forma $ F(x, y) = 0 $.
\end{enumerate}

De exemplu, jumătatea superioară a cercului unitate poate fi dată:
\begin{itemize}
  \item parametric, în forma:
    \[
      \begin{cases}
        x &= \cos t \\
        y &= \sin t
      \end{cases}, t \in (0, \pi);
    \]
  \item cartezian implicit, sub forma: $ y = \sqrt{1 - x^2} $;
  \item cartezian explicit, sub forma $ x^2 + y^2 = 1, y > 0 $.
\end{itemize}

%%%%%%%%%%%%%%%%%%%%%%%%%%%%%%%%%%%%%%%%%%%%%%%%%%%%%%%%%%%%%%%%%%%%%%
\section{Tangenta și normala la curbă}

Dată o curbă plană, în funcție de forma în care a fost prezentată,
putem scrie ecuațiile pentru tangenta și normala la curbă.

\index{curbe plane!tangenta}
\index{curbe plane!normala}
Pentru reprezentarea parametrică:
\begin{itemize}
  \item tangenta în punctul $ P(x(t), y(t)) $ are ecuația:
    \[
      \dfrac{x - x(t)}{x'(t)} = \dfrac{y - y(t)}{y'(t)};
    \]
  \item normala în punctul $ P(x(t), y(t)) $ are ecuația:
    \[
      x'(t)(x - x(t)) + y'(t)(y - y(t)) = 0.
    \]
\end{itemize}

Dacă avem curba dată cu reprezentarea carteziană implicită, presupunînd că
funcția $ F $ este regulată:
\begin{itemize}
  \item tangenta în punctul $ A(x_0, y_0) \in \RR^2 $ are ecuația:
    \[
      (x - x_0) \dfrac{\partial F}{\partial x} \Big|_{(x_0, y_0)} + %
      (y - y_0) \dfrac{\partial F}{\partial y} \Big|_{(x_0, y_0)} = 0.
    \]
  \item normala în punctul $ A(x_0, y_0) \in \RR^2 $ are ecuația:
    \[
      \dfrac{x - x_0}{F_x(x_0, y_0)} = \dfrac{y - y_0}{F_y(x_0, y_0)},
    \]
    unde am notat cu $ F_x = \dfrac{\partial F}{\partial x} $ și
    $ F_y = \dfrac{\partial F}{\partial y} $.
\end{itemize}

%%%%%%%%%%%%%%%%%%%%%%%%%%%%%%%%%%%%%%%%%%%%%%%%%%%%%%%%%%%%%%%%%%%%%%

\section{Curbura unei curbe plane}

\index{curbe plane!curbura}
\index{curbe plane!raza de curbură}
Fie $ \alpha $ o curbă regulată, dată de ecuații parametrice cu $ t \in (a, b) $.

Se definește funcția:
\[
  k : (a, b) \to \RR, \quad k(t) = \dfrac{x'y^{\prime\prime} - x^{\prime\prime} y'}{%
  (x^{\prime 2} + y^{\prime 2})^{3/2}},
\]
care se numește \emph{curbura} curbei $ \alpha $, iar $ \dfrac{1}{|k|} $ se
numește \emph{raza de curbură}.

%%%%%%%%%%%%%%%%%%%%%%%%%%%%%%%%%%%%%%%%%%%%%%%%%%%%%%%%%%%%%%%%%%%%%%

\section{Reperul Frenet în plan}
\index{curbe plane!reperul Frenet}

Dată orice curbă plană, există un sistem de referință (reper sau sistem de
coordonate) care este mobil și se deplasează odată cu curba. Acesta este
dat de \emph{versorul tangent} $\vec{T} $ și \emph{versorul normal} $ \vec{N} $
la curbă, împreună formînd \emph{reperul Frenet.}

Pornind cu o curbă reprezentată parametric $ \alpha = \alpha(x(t), y(t)) $,
versorii de mai sus se definesc astfel:
\begin{align*}
  \vec{T}(t) &= \dfrac{\alpha'(t)}{||\alpha'(t)||} = %
  \dfrac{(x'(t), y'(t))}{\sqrt{x(t)^{\prime 2} + y(t)^{\prime 2}}} \\
  \vec{N}(t) &= R_{\frac{\pi}{2}} \vec{T}(t) = %
  \dfrac{(-y'(t), x'(t))}{||\alpha'(t)||}.
\end{align*}

Astfel, reperul $(\vec{N}, \vec{T}) $ este unul ortonormat și mobil.

%%%%%%%%%%%%%%%%%%%%%%%%%%%%%%%%%%%%%%%%%%%%%%%%%%%%%%%%%%%%%%%%%%%%%%

\section{Cercul osculator}
\index{curbe plane!cercul osculator}

Se definește cercul osculator ca fiind cercul de rază egală cu raza de
curbură a curbei date, iar centrul cercului este în punctul de coordonate
$ (x_0, y_0) $, definite prin ecuațiile:
\begin{align*}
  x_0 &= x - y' \cdot \dfrac{x^{\prime 2} + y^{\prime 2}}{%
  x'y^{\prime\prime} - x^{\prime\prime}y'} \\
    y_0 &= y + x' \cdot \dfrac{x^{\prime 2} + y^{\prime 2}}{%
    x'y^{\prime\prime} - x^{\prime\prime}y'}.
\end{align*}

%%%%%%%%%%%%%%%%%%%%%%%%%%%%%%%%%%%%%%%%%%%%%%%%%%%%%%%%%%%%%%%%%%%%%%

\section{Curbe în spațiu --- Reprezentare}

\index{curbe în spațiu!ecuații parametrice}
\index{curbe în spațiu!ecuații carteziene implicite}

Similar cu cazul plan, curbele spațiale pot fi reprezentate prin:
\begin{enumerate}[(1)]
  \item ecuații parametrice, sub forma:
    \[
      \alpha : \begin{cases}
        x &= x(t) \\
        y &= y(t) \\
        z &= z(t)
      \end{cases}, t \in (a, b);
    \]
  \item ecuații carteziene implicite, sub forma unei intersecții de
    două suprafețe, sub forma generală:
    \[
      \begin{cases}
        f(x, y, z) &= 0 \\
        g(x, y, z) &= 0
      \end{cases}.
    \]
\end{enumerate}


%%%%%%%%%%%%%%%%%%%%%%%%%%%%%%%%%%%%%%%%%%%%%%%%%%%%%%%%%%%%%%%%%%%%%%

\section{Tangenta și planul normal}
\index{curbe în spațiu!tangenta}
\index{curbe în spațiu!planul normal}

În cazul în care curba a fost dată cu reprezentarea parametrică,
$ \alpha = \alpha(x(t), y(t), z(t)), t \in (a, b) $, fie un punct
$ P(x(t), y(t), z(t)) $ pe curbă.

\begin{itemize}
  \item tangenta în punctul $ P $ are ecuația:
    \[
      \dfrac{x - x(t)}{x'(t)} = \dfrac{y - y(t)}{y'(t)} = \dfrac{z - z(t)}{z'(t)};
    \]
  \item planul normal în punctul $ P $ are ecuația:
    \[
      x'(t)(x - x(t)) + y'(t)(y - y(t)) + z'(t)(z - z(t)) = 0.
    \]
\end{itemize}

Dacă avem curba dată cartezian implicit, fie $ P(x_0, y_0, z_0) $ un punct
pe curbă.
\begin{itemize}
  \item tangenta în punctul $ P $ are ecuația:
    \[
      \dfrac{x - x_0}{\frac{D(f, g)}{D(y_0, z_0)}} = %
      \dfrac{y - y_0}{\frac{D(f, g)}{D(z_0, x_0)}} = %
      \dfrac{z - z_0}{\frac{D(f, g)}{D(x_0, y_0)}};
    \]
  \item planul normal în $ P $ are ecuația:
    \[
      (x - x_0) \dfrac{D(f, g)}{D(y_0, z_0)} + (y - y_0) \dfrac{D(f, g)}{D_(z_0, x_0)} + %
      (z - z_0) \dfrac{D(f, g)}{D(x_0, y_0)},
    \]
\end{itemize}

unde am notat în ambele ecuații jacobienii:
\[
  \dfrac{D(f, g)}{D(y_0, z_0)} = \begin{vmatrix}
  \dfrac{\partial f}{\partial y} & \dfrac{\partial f}{\partial z} \\
                                 & \\
  \dfrac{\partial g}{\partial y} & \dfrac{\partial g}{\partial z}
\end{vmatrix}_{(y_0, z_0)}
\]
și similar pentru celelalte coordonate.

%%%%%%%%%%%%%%%%%%%%%%%%%%%%%%%%%%%%%%%%%%%%%%%%%%%%%%%%%%%%%%%%%%%%%%

\section{Elementele Frenet în spațiu}
\index{curbe în spațiu!elemente Frenet}

Dată o curbă reprezentată parametric $ \alpha = \alpha(t) $, se definesc
elementele Frenet:
\begin{enumerate}[(1)]
  \item cîmpul tangent: $ \vec{T} = \dfrac{\alpha'}{||\alpha'||} $;
  \item cîmpul binormal: $ \vec{B} = \dfrac{\alpha' \times \alpha^{\prime\prime}}{%
    \alpha' \times \alpha^{\prime\prime}} $;
  \item cîmpul normal principal: $ \vec{N} = \vec{B} \times \vec{T} $;
  \item curbura: $ k = \dfrac{||\alpha' \times \alpha^{\prime\prime}||}{%
    ||\alpha'||^3} $;
  \item torsiunea: $ \tau = \dfrac{\langle \alpha' \times \alpha^{\prime\prime}, %
    \alpha^{\prime\prime\prime} \rangle}{||\alpha' \times \alpha^{\prime\prime}||^2} $.
\end{enumerate}

În funcție de aceste elemente, distingem următoarele situații:
\begin{itemize}
  \item dacă $ k = 0 $, curba $ \alpha $ este o (porțiune dintr-o) dreaptă;
  \item dacă $ \tau = 0 $, curba $ \alpha $ este o curbă plană;
  \item dacă $ k $ este constantă pozitivă și $ \tau = 0 $, curba $ \alpha $
    este un arc de cerc;
  \item dacă raportul $ \dfrac{\tau}{k} $ este constant, curba $ \alpha $ este
    o (porțiune dintr-o) elice cilindrică;
  \item dacă $ k $ este o constantă pozitivă și $ \tau $ este o constantă nenulă,
    atunci $ \alpha $ este (un arc dintr-o) elice circulară.
\end{itemize}


%%%%%%%%%%%%%%%%%%%%%%%%%%%%%%%%%%%%%%%%%%%%%%%%%%%%%%%%%%%%%%%%%%%%%%

\section{Reperul Frenet în spațiu}
\index{curbe în spațiu!reperul Frenet}

Muchiile reperului Frenet în spațiu într-un punct $ t = t_0 $ sînt date de:
\begin{itemize}
  \item dreapta determinată de punctul $ \alpha(t_0) $ și vectorul tangent $ \vec{T}(t_0) $;
  \item dreapta determinată de punctul $ \alpha(t_0) $ și vectorul normal $ \vec{N}(t_0) $;
  \item dreapta determinată de punctul $ \alpha(t_0) $ și vectorul binormal $ \vec{B}(t_0) $.
\end{itemize}

Fețele reperului Frenet sînt:
\begin{itemize}
  \item planul osculator, care este planul determinat de vectorii $ \vec{T}(t) $ și $ \vec{N}(t) $;
  \item planul normal, detemrinat de vectorii $ \vec{N}(t)$ și $\vec{B}(t) $;
  \item planul rectificant, determinat de $ \vec{T}(t) $ și $ \vec{B}(t) $.
\end{itemize}

\begin{remark}\label{rk:plan-vectori-punct}
  Amintim, în general, ecuația planului determinat de vectorii $ \vec{v} = (a, b, c) $ și
  $ \vec{w} = (m, n, p) $, care trece prin punctul $ P(x_0, y_0, z_0) $ are ecuația
  generală:
  \[
    \begin{vmatrix}
      x - x_0 & y - y_0 & z - z_0 \\
      a & b & c \\
      m & n & p
    \end{vmatrix} = 0.
  \]
\end{remark}

%%%%%%%%%%%%%%%%%%%%%%%%%%%%%%%%%%%%%%%%%%%%%%%%%%%%%%%%%%%%%%%%%%%%%%

\section{Exerciții}

\indent\indent 1. Fie curba plană $ \alpha(t) = (t^2, 3t), t \in \RR $. 
Aflați normala și tangenta la curbă în punctul $ A(1, -3) $.

\vspace{1cm}

2. Fie curba $ \alpha : x^2 - y^3 - 3 = 0 $.
\begin{enumerate}[(a)]
  \item aflați tangenta și normala la curbă în punctul $ A(-2, 1) $;
  \item parametrizați curba $ \alpha $.
\end{enumerate}

\vspace{1cm}

3. Fie parabola $ \alpha(t) = (t, t^2), t \in \RR $.

Aflați ecuația cercului osculator în punctul $ A(-2, 4) $ al curbei și ecuația
carteziană a curbei.

\vspace{1cm}

4. Fie curba $ \alpha(t) = (2 \cos t, 2 \sin t, t), t \in \RR $.
\begin{enumerate}[(a)]
  \item Aflați elementele Frenet în punctul $ A(-2, 0, \pi) $;
  \item Aflați ecuațiile carteziene ale curbei;
  \item Arătați că $ \alpha $ este o elice;
  \item Aflați muchiile și fețele reperului Frenet.
\end{enumerate}

\vspace{1cm}

5. Se dă curba $ \alpha(t) = (t + t^2, t^2 - t, t^2 - t), t \in \RR $.
Arătați că:
\begin{enumerate}[(a)]
  \item $ \alpha $ are planul osculator independent de $ t $ ;
  \item are torsiunea nulă;
  \item are vectorul binormal independent de $ t $.
\end{enumerate}

\vspace{1cm}

6. Fie curba:
\[
  \alpha : \begin{cases}
    x^2 + y^2 + z^2 &= 2  \\
    z &= 1
  \end{cases}.
\]

\begin{enumerate}[(a)]
  \item Aflați dreapta tangentă și planul normal la curbă în $ A(-1, 0, 1) $;
  \item Determinați o parametrizare a curbei.
\end{enumerate}

